% CPSY 1291 -- Recitation handout 13
% Compile:  latexmk -lualatex handout-13-project-clinic-i.tex
\documentclass[11pt]{article}
\usepackage{cpsy1291-handout}

\begin{document}

\handouthead{13}{Project clinic I}
  {A question you can answer, data you can actually get, and a baseline by Friday}
  {Week 13 \ $\cdot$\ project work begins Thu 12/3; presentations Mon 12/21}

\section*{What this session is for}
\addcontentsline{toc}{section}{What this session is for}

You have about two and a half weeks and roughly 31 hours each. The single
largest risk to your project is not that the model fails --- a negative result,
clearly analyzed, is a perfectly good project and the syllabus says so. The risk
is spending twelve days on setup and having nothing to say on the thirteenth.

So this handout is about the first three days: choosing a question that is
answerable at this scale, getting data that actually loads, and having a working
end-to-end result --- however bad --- before you improve anything.

\begin{keybox}
\ask{The one rule.} \textbf{Get to a bad answer fast.} A pipeline that runs
end to end on ten items and produces a wrong number is worth more on day three
than a beautiful half-finished component, because it tells you what the real
obstacles are. Everything after that is improvement, and improvement is
schedulable in a way that discovery is not.
\end{keybox}

\section{Scoping a question}

\subsection{The three shapes, and what each needs by day three}

\begin{center}
\begin{tabular}{@{}p{0.28\textwidth}p{0.32\textwidth}p{0.32\textwidth}@{}}
\toprule
Shape & The risk & What you must have by day 3 \\
\midrule
\textbf{Extend a paper} & the original result does not reproduce & the original result, reproduced on a subset \\
\textbf{Improve a model} & no measurable axis & a baseline number on your metric \\
\textbf{Compare with data} & the data never arrives & the data loaded and plotted \\
\bottomrule
\end{tabular}
\end{center}

Every one of these is a \emph{de-risking} step rather than a result. Do it first,
and if it fails you still have eleven days to change course.

\subsection{The sentence test}

Write your question in this form, filling every blank with something concrete:

\begin{keybox}
\ask{``We measure \underline{\hspace{3.2cm}} on \underline{\hspace{3.2cm}},
comparing \underline{\hspace{2.4cm}} against \underline{\hspace{2.4cm}}, and we
expect \underline{\hspace{3.2cm}}.''}
\end{keybox}

If any blank is vague, the project is not yet scoped, and no amount of coding
will fix that. Two failure patterns to check for:

\begin{itemize}
  \item \textbf{No comparison.} ``We will see whether a CNN represents shape'' has
        no second condition and therefore no possible result. ``We compare a
        CNN's accuracy on shape-preserving versus texture-preserving distortions''
        does.
  \item \textbf{No expectation.} If you cannot say what you expect, you will not
        be able to tell an interesting result from a bug. Write the expectation
        down \emph{before} you run anything, and keep it --- being wrong is a
        finding, and being wrong without a record is just confusion.
\end{itemize}

\subsection{Scope: what fits in 31 hours}

\begin{center}
\begin{tabular}{@{}ll@{}}
\toprule
Fits & Does not fit \\
\midrule
fine-tuning a pretrained model & training a large model from scratch \\
a few thousand images & ImageNet \\
2--4 conditions & a factorial sweep \\
one dataset, done carefully & three datasets, done carelessly \\
a laptop or Colab's free tier & anything needing a reserved GPU \\
\bottomrule
\end{tabular}
\end{center}

\begin{keybox}
\ask{Assume the network is slow.} Downloads in this course have run at
kilobytes per second. Start every download on day one, in the background, and
build against a small local subset while it runs. A project blocked on a
download at day ten is the most avoidable failure there is.
\end{keybox}

\section{Getting data}

\subsection{Where to look, in order}

\begin{enumerate}
  \item \textbf{Datasets you already have.} Assignments 1--5 shipped stimulus
        sets, activation matrices, and real neural data. Reusing one is not
        cheating; it is good sense, and it means day one is a loading step
        rather than a search.
  \item \textbf{The paper's own release.} Many papers post data on OSF, Zenodo,
        or GitHub. Check for a ``Data availability'' statement before you write
        to anyone.
  \item \textbf{Standard repositories} --- OpenNeuro for neuroimaging, CRCNS for
        electrophysiology, the Natural Scenes Dataset, Brain-Score's public
        benchmarks, HuggingFace Datasets for text and images.
  \item \textbf{Ask.} Authors often share on request, but budget a week for a
        reply and never make the project depend on one arriving.
\end{enumerate}

\subsection{Load it, then look at it}

Before any analysis, and without exception:

\begin{lstlisting}
print(X.shape, X.dtype)
print(np.isnan(X).sum(), np.isinf(X).sum())
print(X.min(), X.max(), X.mean())
print(np.unique(y, return_counts=True))     # class balance
plt.imshow(X[0]); plt.show()                # LOOK at one item
\end{lstlisting}

Five lines, and they catch: transposed dimensions, images in $[0,255]$ where the
model expects $[0,1]$, a class that occurs eleven times, dead units that will
produce \texttt{nan} correlations, and files that decoded into noise. Every one
of those has cost somebody a week.

\subsection{Write a loader, once}

\begin{lstlisting}
def load_data(root, subset=None):
    """Return X, y, meta. The ONLY place that touches the filesystem."""
    ...
    return X, y, meta
\end{lstlisting}

Everything downstream calls this. When you discover on day nine that the labels
were off by one, you fix it in one place and rerun --- rather than hunting for
six notebooks that each did their own loading slightly differently.

\section{The end-to-end skeleton}

Write this on day two, with every step deliberately stupid, and run it on ten
items.

\begin{lstlisting}
X, y, meta = load_data(ROOT, subset=10)     # 1. ten items
model      = get_model()                    # 2. pretrained, off the shelf
preds      = predict(model, X)              # 3. no batching, no speed
score      = evaluate(preds, y)             # 4. one number
plot(score)                                 # 5. one figure
\end{lstlisting}

\begin{keybox}
\ask{Why ten items.} Because it runs in seconds, so you can iterate; because
you can verify every number by hand; and because scaling from 10 to 10,000 is a
known, boring change, whereas scaling from ``nothing runs'' to ``something runs''
is not. When the skeleton works, raise \texttt{subset} and go and have lunch.
\end{keybox}

\subsection{Two baselines you always need}

\textbf{Chance.} What score does a model that knows nothing get? For a classifier
with imbalanced classes it is the majority proportion, not $1/k$. Compute it, do
not assume it.

\textbf{A trivial model.} Whatever the sophisticated thing is, run the dumb thing
too --- raw pixels instead of network features, logistic regression instead of a
fine-tuned transformer, mean luminance instead of a shape descriptor. If the dumb
thing does nearly as well, you have learned the most important fact about your
problem, and you have learned it on day three rather than in the last slide of
your talk.

\section{Reproducibility, cheaply}

Four habits, each costing under a minute, that will save your group hours:

\begin{enumerate}
  \item \textbf{One seed, set once, at the top.} \texttt{np.random.default\_rng(0)},
        \texttt{torch.manual\_seed(0)}. And remember Recitation 3: one seed makes
        it reproducible, not meaningful. For any headline number, run several.
  \item \textbf{Save results, not just figures.} Write the numbers to a
        \texttt{.csv} or \texttt{.npz} as you go. Regenerating a figure from a
        saved array takes ten seconds; regenerating it from scratch takes an
        hour and may not match.
  \item \textbf{Never edit a figure by hand.} If it needs a change, change the
        code. Hand-edited figures cannot be regenerated, and you will need to
        regenerate them.
  \item \textbf{One notebook per question, not one per person.} Three people
        working in three copies of the same notebook is the most common way a
        group project loses a day to merge conflicts. Split by \emph{question},
        with a shared \texttt{utils.py} that all three import.
\end{enumerate}

\section{A schedule that works}

\begin{center}
\begin{tabular}{@{}ll@{}}
\toprule
By & You have \\
\midrule
day 3  & data loaded and looked at; the skeleton runs on 10 items; chance and a trivial baseline \\
day 6  & the real comparison run once, at full size, with one figure \\
day 9  & the comparison run across seeds, with error bars; one control \\
day 12 & every figure final; the result stated in one sentence \\
day 14 & slides built and rehearsed once, out loud, with a timer \\
\bottomrule
\end{tabular}
\end{center}

\begin{keybox}
\ask{The deadline that matters is day 9, not day 14.} If the main comparison is
not running by day 9 you do not have time to add error bars, and a result without
error bars cannot be defended --- which is the single most common way a
technically fine project loses marks.
\end{keybox}

\section{Ten questions to ask your own project}

\begin{enumerate}
  \item What is the comparison? If there is only one condition, there is no result.
  \item What did you \emph{expect}, and did you write it down first?
  \item What is chance, and did you compute it or assume it?
  \item What does the trivial baseline get?
  \item How many seeds, and what is the spread?
  \item What is the noise ceiling, if you are comparing against data?
  \item What single change would most likely overturn your conclusion?
  \item Did you touch the test set more than once?
  \item Can someone else run your pipeline from your loader?
  \item If the result had come out the other way, would the project still work?
\end{enumerate}

Question 10 is the one to ask on day one. If the answer is no, the project is a
bet rather than an experiment, and you have two weeks.

\end{document}
